\scanpage{205}\mspage{154}
\[
X_{0/x_0}\longrightarrow X_{1/x_1},\qquad X_{0/x_0}\mathbin{\unsure{-|\to}} X_{2/x_2}
\]

\inserted{(définis à iso. uniques près)} ne sont autres que \unsure{les|\deleted{les}} localisées $\alpha_{1/x_0}$ et $\alpha_{2/x_0}$ de $\alpha_1:X_0\mathbin{-}X_1$ et $\alpha_2:X_0\mathbin{-}X_2$. Cela montre donc qu'on a un iso.

\diagram{\(\unsure{\mathscr X|\mathcal X}_{/x_0}\simeq\int \begin{tikzcd}[ampersand replacement=\&,row sep=0pt,column sep=1.2em,cramped] \& X_{1/x_1} \\ X_{0/x_0}\arrow[ur]\arrow[dr] \& \\ \& X_{2/x_2} \end{tikzcd},\)}

et l'on admettra aussi que la flèche $(\varphi_Q)_{/x_0}:\mathscr X_{/x_0}\mathbin{-}Y_{/y}$ s'identifie à

\[
\varphi(Q_{/x_0}).
\]

\unsure{Pour} exprimer la \inserted{$W$-}lissité de $\varphi_Q$ en $x_0$\inserted{$\in\mathrm{Ob}\,X_0\subset\mathrm{Ob}\,\mathscr X$}, on est donc ramené à exprimer la $W$-asphéricité des fibres de

\diagram{\(\varphi(Q_{/x_0}):\ \int \begin{tikzcd}[ampersand replacement=\&,row sep=0pt,column sep=1.2em,cramped] \& X_{1/x_1} \\ X_{0/x_0}\arrow[ur]\arrow[dr] \& \\ \& X_{2/x_2} \end{tikzcd}\longrightarrow Y_{/y}.\)}

Pour expliciter les fibres, on a besoin d'un autre lemme auxiliaire

\uline{Lemme 2} $^{\star}$ Soit un carré commutatif dans $\mathrm{Cat}$

\margin{NB \unsure{C'est le} b) \unsure{du Th} 4 \deleted{\unsure{page 148}}} % slanted

\diagram{\(Q=\biggl(\begin{tikzcd}[ampersand replacement=\&,row sep=1.0em,column sep=1.8em] X_0\arrow[r,"\alpha_1"]\arrow[d,"\alpha_2"'] \& X_1\arrow[d,"f_1"] \\ X_2\arrow[r,"\deleted{\unsure{f_0}}" above,"f_2" below] \& Y \end{tikzcd}\biggr),\)}

et \deleted{\unsure{soit}}

\diagram{\(\varphi_Q:\bigl(\int \begin{tikzcd}[ampersand replacement=\&,row sep=0pt,column sep=1.2em,cramped] \& X_1 \\ X_0\arrow[ur]\arrow[dr,dash] \& \\ \& X_2 \end{tikzcd}\bigr)\longrightarrow Y,\)}
% archive sheet number: 205
